\section{Introduction}\label{sec:meta_prior_introduction}

Simulation optimization becomes difficult when a policy is written as a long
ordered profile and each evaluation returns both a noisy objective and a
probabilistic safety outcome. The profile may contain many coordinates, but the
coordinates often arise from a smaller collection of response shapes. A new
task therefore contains two kinds of uncertainty: which shapes matter in the
target system, and how the target system changes their mean response and
accumulated risk.

Historical simulations can guide both structural initialization and a
source-to-target function prior. The project retains a frozen source-informed
combined frontend that selects initial policies before target observations; its completed
component evidence is reported in Section~\ref{sec:mp_legacy_frontend}.
A structural meta-prior learns
which functions explain source responses, assigns a distribution to their
coefficients and residual risk, and lets paid target observations update that
distribution. These components address different decisions: the frontend chooses
where to start, while the prior describes the response and risk functions
updated during target search. This distinction matters in chance-constrained problems:
source evidence may improve early prediction, yet a mismatched safety model
can move the recommended policy in the wrong direction. The method must
therefore preserve target correction, negative transfer, and independent
verification in the same protocol.

We study this bounded problem in a declared public representation. Each
policy is mapped to an eight-dimensional coordinate and evaluated through a
finite dictionary of 68 centered and scaled functions. Separate source-task
models describe the objective mean, the constraint mean, and the logarithm of
cumulative risk variance. The target is evaluated with source-selected, same-source
full-dictionary, and target-only controls. Target observations update the
function posterior, while physical-risk family weights are updated from risk
labels. A
correlated Gaussian knowledge-gradient rule then values a candidate by its
expected verified deployment utility under the remaining evaluation and
verification budget.

The resulting information contract is deliberately narrower than a general
representation-learning claim. The coordinate is declared rather than
learned from trajectories; the dictionary is finite and not orthogonalized;
and the risk head models total log variance rather than recovering separate
physical common and local factors. These choices make the source-to-target
path auditable and allow negative transfer to remain visible.

The paper makes three contributions. First, it specifies a reproducible
source-to-target structural prior in which source observations learn function
support and cross-task coefficient variation, while target observations can
remove a mismatched structure. Second, it couples this posterior to a
verification-aware decision rule and records the distinct roles of prediction,
candidate ranking, and independent certification. Third, it evaluates the
contract under paired source archives, initial designs, target budgets, and
verification streams. The results show an early-budget advantage under
matched mechanisms, a small and unstable contribution from learned risk
structure, and no universal advantage for the knowledge-gradient or for
operational deployment.

This separation gives a precise interpretation of the evidence. A reduction
in prediction error is not treated as a deployment certificate; a small
knowledge-gradient increment is not treated as proof of a general sampling
advantage; and an infeasible target task is retained as a failure of the
end-to-end deployment protocol. The conclusions therefore apply to the
declared representation and finite dictionary, while trajectory encoders,
factor-level risk recovery, and general high-dimensional optimization remain
open problems.
